\section{Black and white}
The \robot\ can colour in its drawings.
We will start with black and white colouring. The command {\idt paint()}
lets black paint escape from the back of the \robot. The paint
spreads over the largest possible part of the screen bounded by a black line.
An example is given in figure \ref{fig:peinyy} below.
\begin{figure}[htbp]
 \includevspacePostScript[2.9cm]{PeintureYinYang}{0.95}{0.95}
 \caption{Steps in painting a drawing}\label{fig:peinyy}
\end{figure}
We start by drawing the outline of the area to be painted, then we place the \robot\ in the area to be
painted or on its edge
so
that the back of the \robot, through which the paint escapes, is in the area to be painted.
The paint spreads line by line above,
to the left of, to the right of and below the back point of the \robot. The area covered is the largest one
containing
the back of the
\robot\ and
bounded by a continuous black line. The \robot\ can paint without moving, whether its pen is up or down.
\needspace{20\baselineskip}
\begin{exemple}[Yin Yang]{\rm
 The program below draws the Chinese Yin Yang symbol:

 % figure et programme insecables (a un saut de page, la figure disparaissait)
 \noindent\begin{minipage}{\linewidth}
 \vbox to 0pt{\vspace*{25mm}\includePostScript{YinYang}{0.95}{0.95}}
 \input python/YinYang.tex
 \end{minipage}
\noindent $\Box$ }\end{exemple}

A frequent mistake is to leave a hole in the boundary of the area to be painted through
which the paint will leak. A hole of a single pixel is enough to cause a leak. In case of a leak, we can
stop the painting by pressing one of the keys
\begin{picture}(12,12)(0, 2) %touche d'echappement
  \put(0,0){\framebox(12,12)[lb]{}}
  \put(1.5,2){\makebox(0,0)[lb]{\tiny esc}}
  \end{picture}\hspace{0.3ex} or
  \begin{picture}(12,12)(0, 2) %touche §
  \put(0,0){\framebox(12,12)[lb]{}}
  \put(1.5,2){\makebox(0,0)[lb]{\tiny \S}}
  \end{picture}\hspace{0.3ex} or
  \begin{picture}(12,12)(0, 2) %touche #
  \put(0,0){\framebox(12,12)[lb]{}}
  \put(1.5,2){\makebox(0,0)[lb]{\tiny \#}}
  \end{picture}\hspace{0.3ex}. The \robot\ goes on drawing.
  Press one of these keys at least a second time to stop it for good.

\section{Colour}

To draw and paint, the \robot\ has pens in the colours white, yellow, blue, violet, green,
red, indigo and black. The command {\idt pc} (\underline{p}en \underline{c}olour) chooses the colour of the pen, which is given as a parameter. For example, after the command {\idt pc(red)}
all lines and painting are in red. The command {\idt pc(black)} returns
to black lines and painting. The colour names {\idt white}, {\idt yellow}, {\idt blue},
{\idt violet}, {\idt green}, {\idt red}, {\idt indigo} and {\idt black} are constants
declared by {\idt LibRobot}, which stand for the colour numbers $0$, $1$, $2$, \ldots, $7$
(in this order): {\idt pc(red)} is therefore the same thing as {\idt pc(5)}, but much more
readable. As we saw in chapter 9, the command {\idt pw(n)} chooses
the thickness of the pen; it is independent of its colour.

On a colour screen, the \robot\ draws and paints with the colours of the first
line below.

When we do not have a colour screen but a screen with several levels of grey,
each pixel of which can be more or less dark, the
colours are represented by greys, in the order white,
yellow, blue, violet, green, red, indigo and black from lightest to darkest.

When we only have a black and white screen, the colours are represented by
grey shades made of patterns of white dots and black dots side by side. The order is
again white,
yellow, blue, violet, green, red, indigo and black from lightest to darkest. The darker the shades,
the denser the small dots used to represent them on the screen, as
can be seen on the second line below:\\[1.8cm]
\includePostScript{CouleursCouleur}{0.95}{0.95}\ \\[1.8cm]
\includePostScript{Couleurs}{0.95}{0.95}\ \\

When painting, the edges of the pattern must be drawn with no possible leak, in any colour except white.

\needspace{30\baselineskip} % la fleur et le debut du programme sur la meme page
\begin{exemple}[Flower in colour]{\rm
 The program below paints the flower opposite in colour:\\[5cm]
 \includePostScript{FLeurStyliseeCouleur}{0.95}{0.95}%
 \vspace*{-5cm}
 \input python/StylizedFlowerColour.tex
\noindent $\Box$ }\end{exemple}

\needspace{6\baselineskip}\noindent \rule{137mm}{0.1mm}\par\nopagebreak %---------------------------------------------
\begin{exercices}{Write \python\ programs to draw and colour in black and white:

\tiret the tilings of the square of figure \ref{fig:pavac} below;

\tiret the mosaics of the Hall of the Abencerrages (14th century) of the Al\-ham\-bra,
the residence of the Moorish kings in Granada, Andalusia, as they were copied
by the artist Escher in 1936 \cite{escher1} (figure \ref{fig:mosAlham} on page \pageref{fig:mosAlham}).

\noindent $\Box$}\end{exercices}
\begin{figure}[htbp]
 \includevspacePostScript[4.3cm]{Pavages.DuCarre}{0.95}{0.95}
 \caption{Tilings of the square}\label{fig:pavac}
\end{figure}
\begin{figure}[htbp]
 \includevspacePostScript[4.6cm]{MosaiqueAlhambra123}{0.95}{0.95}
 \caption{Mosaics of the Alhambra}\label{fig:mosAlham}
\end{figure}
\begin{exercices}{Write \python\ programs to draw and colour in col\-our:

\tiret the ball and the three-colour tiling of figure \ref{fig:balpav} (page \pageref{fig:balpav});

\tiret the mosaic of the Al\-ham\-bra \cite{escher1} of figure \ref{fig:mosai4}
        (page \pageref{fig:mosai4});

\tiret the tiling of disjoint squares of figure \ref{fig:pdcc} (page \pageref{fig:pdcc});

\tiret the tilings of figures
\ref{fig:PVdV} (page \pageref{fig:PVdV}),
\ref{fig:PVdV1l} (page \pageref{fig:PVdV1l}),
\ref{fig:PVdV2l} (page \pageref{fig:PVdV2l})
and
\ref{fig:PVdV4l} (page \pageref{fig:PVdV4l})
using the ornamental motif by N. J. van de Vecht as it is
painted in figure \ref{fig:mVdV} (page \pageref{fig:mVdV}).

\noindent $\Box$}\end{exercices}
\begin{figure}[htbp]
\includevspacePostScript[4.4cm]{BallePavageCouleur}{0.95}{0.95}
 \caption{Ball, three-colour tiling}\label{fig:balpav}
\end{figure}
\vfill
\begin{figure}[htbp]
 \includevspacePostScript[5.6cm]{MosaiqueAlhambra4}{0.95}{0.95}
 \caption{Mosaic of the Alhambra in colour}\label{fig:mosai4}
\end{figure}
\begin{figure}[htbp]
\vspace*{\textwidth}%
\vspace*{25mm}%
\includevspacePostScript[8.8cm]{PavageDis.Car.Cou}{0.95}{0.95}%
\vspace*{-\textwidth}%
\vspace*{-25mm}%
\caption{Tiling of disjoint squares in colour}\label{fig:pdcc}
\end{figure}
\begin{exercice}{What is the motif used to make the black and white tiling of figure \ref{fig:pavhexond}
(page \pageref{fig:pavhexond})?
\noindent $\Box$}\end{exercice}
\begin{figure}[htbp]
 \includevspacePostScript[13.7cm]{PavageHexagonalOndule}{0.95}{0.95}
 \caption{Tiling of the plane in black and white}\label{fig:pavhexond}
\end{figure}
%%%%%%%%%%%%%%%%%%%%%%%%%%%%%%%%%%%%%%%%%%%%%%%%%%%%%%%%
\begin{corrige}{
We give one program; the others are in the folder {\idt book/\allowbreak 12\_painting}:
\input python/AlhambraMosaic3.tex
\noindent $\Box$}\end{corrige}
\begin{corrige}{
The motif of the black and white tiling of the plane of figure \ref{fig:pavhexond}
(page \pageref{fig:pavhexond}) is given in figure \ref{fig:motpavhexond}
(page \pageref{fig:motpavhexond}).
\noindent $\Box$}\end{corrige}
\begin{figure}[htbp]
 \includevspacePostScript[7.4cm]{MotifDuPavageHexagonalOndule}{0.95}{0.95}
 \caption{Motif of the tiling of figure \protect{\ref{fig:pavhexond}}}\label{fig:motpavhexond}
\end{figure}
\begin{corrige}{
We give one program; the others are in the folder {\idt book/\allowbreak 12\_painting}:
\label{prog:BalleEnCouleur}
\input python/ColourBall.tex
\noindent $\Box$}\end{corrige}

\vfill

%%%%%%%%%%%%%%%%%%%%%%%%%%%%%%%%%%%%%%%%%%%%%%%%%%%%%%%%%%%%%%%%%%%%%%%%%%%
\noindent \framebox[\textwidth]{
\begin{minipage}{\myparboxwidth}
\setlength{\parindent}{2em}
\small \sf
\vspace{\abovedisplayskip}
\section*{Addenda}

We can imagine that, to paint, the \robot\ places a seed on the screen, behind
its stopping point. If a seed is placed on a pixel of the screen that is already coloured, then
this seed disappears. In particular, if the initial seed placed by the \robot\ is on a non-white pixel,
the painting stops immediately.
If the seed is placed on a white pixel of the screen, then the seed disappears, the white pixel on
which it was is painted in the colour of the
pen and four new seeds are placed to the east, to the west, to the north and to the south of the pixel that
has just been painted. The painting is finished when all the seeds have disappeared.

This painting algorithm is therefore based on the physical representation of the drawing on the screen and
not on the logical representation which would be obtained by recording all the commands
executed by the \robot. As a result, leaks sometimes occur that cannot be explained logically but
that are due to the difference between the logical and physical representations.

Let us take the example of a circle. Mathematically a circle is a continuous line made up
of an infinite set of dimensionless points. On a computer screen the circle will be
represented by a discontinuous line made up of a finite set of ``points'' the size
of a pixel. The mathematical and computer notions of circle are therefore completely
different.

\vspace{\abovedisplayskip}
\end{minipage}
}

%%%%%%%%%%%%%%%%%%%%%%%%%%%%%%%%%%%%%%%%%%%%%%%%%%%%%%%%%%%%%%%%%%%%%%%%%%%%
%%%%%%%%%%%%%%%%%%%%%%%%%%%%%%%%%%%%%%%%%%%%%%%%%%%%%%%%%%%%%%%%%%%%%%%%%%%
\noindent \framebox[\textwidth]{
\begin{minipage}{\myparboxwidth}
\setlength{\parindent}{2em}
\small \sf
\vspace{\abovedisplayskip}
 The approximation is acceptable when the size of the pixels is very small compared with
the diameter of the circle, in the same way as a pair of compasses draws circles that
give a good picture of the mathematical object.
However, for very small circles, we
get some aberrations, visible in the figure below, where the circles shown
in the small frame at the top left of the figure have been enlarged eight times.
Note in particular that the very small circles are square!

\includehvspacePostScript[9.8cm,5.2cm]{LoupeCercles}{0.75}{0.75}

\vspace*{1mm}
On screens with rectangular pixels, circles look
like ellipses. The program must therefore correct the distortion by
logically drawing ellipses in order to obtain circles on
the screen. Circles are closed in the sense that every pixel painted on
the circle has at least one painted neighbouring pixel in one of the eight
directions of the compass rose. This is not the case for the
very elongated ellipses that can be seen in the figure below:
\vspace*{1mm}

\includehvspacePostScript[9.7cm,4.0cm]{LoupeEllipses}{0.75}{0.75}

The holes in the boundary of the elongated ellipses are points where the paint leaks
when we want to paint the inside of the ellipse. To avoid these leaks one would have to use
more complex colouring algorithms based on the logical rather than the physical structure of the
drawing.
\vspace{\abovedisplayskip}
\end{minipage}
}
%%%%%%%%%%%%%%%%%%%%%%%%%%%%%%%%%%%%%%%%%%%%%%%%%%%%%%%%%%%%%%%%%%%%%%%%%%%%
